% !TeX program = xelatex
\documentclass[11pt,a4paper]{article}
\usepackage[margin=23mm]{geometry}
\usepackage{fontspec}
\setmainfont{Latin Modern Roman}
\usepackage{amsmath,amssymb,mathtools}
\usepackage{xcolor,enumitem,booktabs,longtable,array,xurl,hyperref}
\definecolor{muted}{HTML}{53636C}
\hypersetup{colorlinks=true,urlcolor=blue,linkcolor=blue}
\setlength{\parindent}{0pt}
\setlength{\parskip}{5pt}
\newcommand{\R}{\mathbb R}
\newcommand{\E}{\mathbb E}
\newcommand{\N}{\mathbb N}
\newcommand{\ind}{\mathbf 1}
\newcommand{\Var}{\operatorname{Var}}
\newcommand{\Cov}{\operatorname{Cov}}
\newcommand{\supp}{\operatorname{supp}}
\DeclareMathOperator*{\argmin}{arg\,min}
\DeclareMathOperator*{\argmax}{arg\,max}
\newcommand{\entrymeta}[3]{{\sffamily\small\color{muted}\textbf{\texttt{#1}}\quad|\quad #2\par #3\par}}
\newcommand{\Problemref}[1]{\texttt{#1}}
\newcommand{\JELref}[2]{\texttt{#2} (JEL \texttt{#1})}
\begin{document}
\section*{Combining instruments, sign restrictions, and higher moments}
\textbf{Problem ID:} \texttt{C32-139}\quad \textbf{JEL:} \texttt{C32}\par
% BEGIN ECONOMICS STATEMENT
\label{problem:OP-0605}


\label{audited:ECON-021}

{\sffamily\small\color{muted}\textbf{ECON-021} | R | Source-backed agenda; exact formalization is editorial\par}

\subsection*{Definitions and formal target}

This is a parameterized research family. An instance supplies a structural model class \(\mathcal M\) on declared observable and latent process spaces and a time index set. Each \(m\in\mathcal M\) specifies the joint process law, its observable marginal \(P_m\), and \(u_t=B_m\varepsilon_t\), where \(u_t,\varepsilon_t\in\R^d\), \(d\geq1\), and \(B_m\in\R^{d\times d}\) is nonsingular. Declare which variables are observed, how innovations are obtained from the observable process, the reduced-form dynamics and propagation matrices \(A_h(m)\) for declared horizons \(h\in H\subseteq\{0,1,\ldots\}\), and the nuisance domain \(\eta_m\in\mathcal E\). The nuisance description must retain any latent-law information needed for observational compatibility; it need not be finite-dimensional. Any representation \(m=(B,\eta)\) must specify how these objects determine a model and its observable law.

The instance also supplies an explicit inventory of restriction maps, their domains and applicability conditions. This includes shock means, scale normalizations, covariance and cross-time restrictions; whether independence or only specified moment restrictions are imposed; and admissible shock signs, labels and permutations. For a proxy \(Z_t\in\R^k\), declare \(k\), target indices, centering conventions, exclusion moments such as \(\mathbb E_m[Z_t\varepsilon_{jt}]=0\) for non-target shocks, and the precise relevance condition. Declare whether each restriction is unconditional, conditional on an observed regime, or uses a specified latent-regime model. Regime restrictions must state which parameters, including \(B\), are shared. Specify the cumulant orders, all required moment-existence assumptions, and each sign-restricted coordinate of \(A_h(m)B_m\), its horizon and its weak or strict direction. No independence, normalization, regime specification or positive lower bound on identifying strength is implicit. Boundary cases and exact nonidentification are included or excluded explicitly by the instance.

For an observed law \(P\) on the declared observable process space, let
\[
 \mathcal M(P)=\{m\in\mathcal M:P_m=P\},\qquad
 \Theta(P;\mathcal M)=\{\theta(m):m\in\mathcal M(P)\}.
\]
The instance declares a measurable target space \(\mathcal T\) and response map \(\theta:\mathcal M\to\mathcal T\), including response coordinates, horizons, shock units and labeling. For impulse responses the map uses the relevant entries of \(A_h(m)B_m\). On \(\mathcal M(P)\), the shorthand \(\psi_P(B,\eta)\) may denote this map; writing only \(\psi(B)\) requires all other target inputs to be fixed by \(P\). If models are quotiented by admissible sign or permutation equivalences, the target must be invariant on each equivalence class; otherwise retain all permitted labeled responses in the set.

Characterize \(\Theta(P;\mathcal M)\) using the combined identifying information. A sharp characterization must give necessary and sufficient conditions for target membership, including realization by a structural model with the entire observable law \(P\), rather than merely restating the projection above. It need not be an algorithm or finite algebraic test. Restriction maps \(g(P,B,\eta)=0\) and \(q(P,B,\eta)\geq0\), interpreted componentwise over declared index sets, can represent the set as
\[
 \Theta(P;\mathcal M)=\{\psi_P(B,\eta):m=(B,\eta)\in\mathcal M,\ P_m=P,\ g(P,B,\eta)=0,\ q(P,B,\eta)\geq0\}
\]
when they encode maintained restrictions of \(\mathcal M\). Dropping \(P_m=P\) gives only a candidate outer set unless necessity and target-level realizability are established. An empty compatible class gives an empty identified set.

For inference an instance additionally supplies observed data \(X^{(n)}\) on a declared sample space, its sampling law \(Q_{n,m}\), a confidence level \(1-\alpha\) with \(0<\alpha<1\), and nonempty uniformity classes \(\mathcal M_n\). These classes specify dependence, moment conditions, nuisance estimation and admissible weak-identification sequences and restriction failures; no independence of fitted innovations is implicit. For measurable confidence regions \(C_n(X^{(n)})\subseteq\mathcal T\), the target is coverage of each actual structural response \(\theta(m)\). Declare either the finite-sample requirement
\[
 \inf_{m\in\mathcal M_n}Q_{n,m}\{\theta(m)\in C_n(X^{(n)})\}\geq1-\alpha
 \quad\text{for every }n,
\]
or the asymptotic requirement
\[
 \liminf_{n\to\infty}\inf_{m\in\mathcal M_n}
 Q_{n,m}\{\theta(m)\in C_n(X^{(n)})\}\geq1-\alpha.
\]
These do not require simultaneous coverage of the entire identified set. Any additional precision, efficiency or computational requirement must be stated separately with its domain and success criterion; coverage alone permits \(C_n=\mathcal T\).

\subsection*{Clarifications and required assumptions}

Orthogonality to other shocks does not identify a target without relevance and labeling. For validity uncertainty, the instance declares indispensable baseline assumptions defining \(\mathcal M_0\), an indexed inventory \(I\) of uncertain restrictions, and a nonempty collection \(\mathcal V\subseteq2^I\) of permissible retained subsets. Put \(\mathcal M_v=\{m\in\mathcal M_0:\text{all restrictions indexed by }v\text{ hold}\}\) and use \(\mathcal M=\bigcup_{v\in\mathcal V}\mathcal M_v\); restrictions outside a retained subset are not imposed and need not fail. The sharp set is then \(\bigcup_{v\in\mathcal V}\Theta(P;\mathcal M_v)\). Claims of coverage require that every true model in the declared uniformity class belongs to this union; a collection of subsets alone does not ensure this. Known-valid restrictions instead are jointly imposed. For local misspecification, declare the residual maps, permitted deviation sets, their norm or other magnitude convention and sample-size dependence, and construct the corresponding classes \(\mathcal M_n\) with those relaxed restrictions. Exact equalities and inequalities apply only to restrictions still maintained exactly. The target remains the actual model's declared structural response, rather than an unstated pseudo-true parameter. An empty intersection diagnoses incompatibility of the declared restrictions, rather than validating a preferred restriction. All instance inputs above must be supplied before evaluating an application; they are not conventions inferred from the cited agenda.

\subsection*{Variants and changes of scope}

\begin{itemize}
\item Known-valid restrictions give an intersection; unknown-valid subsets give a different, usually wider union.
\item Point identification, informative partial identification, and uniform coverage are separate goals. The whole parameter space achieves coverage without precision.
\item Finite-sample and asymptotic validity differ; dependence and estimated nuisance quantities enter the sampling model.
\end{itemize}

\subsection*{Audit conclusion and source relationship}

[S1], Sections 6.3 and 7.3, supports combining identification information and developing inference. This set and validity-uncertainty criterion are editorial, not a verbatim conjecture.

\subsection*{References and exact source roles}

{\footnotesize\raggedright
\par\noindent\textbf{[S1]} Daniel J. Lewis (2025). \emph{Identification Based on Higher Moments in Macroeconometrics}. Annual Review of Economics 17, 665--693. DOI: 10.1146/annurev-economics-070124-051419. Checked published PDF in the author\textquotesingle s institutional repository.\par \textit{Verified locator / source role:} Sections 6.3 and 7.3. These passages support the methodological agenda; the precise inferential criterion below is editorial.\par \url{https://discovery.ucl.ac.uk/id/eprint/10208055/1/Lewis_annurev-economics-070124-051419.pdf}\par\smallskip
}

\subsection*{Common conventions: Additional-source status and mathematical conventions}

\label{audited:conventions}

Of the 100 supplied ECON entries, 65 distinct problems appear here; 32 close
matches refine earlier entries, and three duplicate questions map to existing
entries. Appendix~\ref{app:audited-crosswalk} lists every destination.
The attachment's P/R/S/U distinctions and source audits are retained. Its
precise mathematical formalizations are editorial unless the individual audit
says otherwise. Candidate subsidy targets ECON-004--006 retain their U flags;
the resolved approval-core question ECON-007 updates OP-0202 above.
The following supplied conventions govern both this part and the attributed
ECON refinements elsewhere in the catalogue, subject to local restrictions.

\textbf{Parameterized questions.} A problem family lists inputs that must be fixed before an instance is evaluated: population, horizon, policies, information, data law, model class and welfare weights. These are required choices, not quantities the citations are assumed to specify. Undefined applications should not be treated as identified estimands.

\textbf{Indices and integrability.} Sets of agents, firms and regions are finite unless a population measure is explicitly used. \(i,j\) index units, \(r\) regions, \(t\) dates and \(h\) horizons. \(\mathbb E\) and \(\Pr\) refer to the declared population and law; \(\mathbf1\{A\}\) is an indicator. Every displayed expectation and discounted sum needs existence in the stated domain. Infinite horizons use a discount in \((0,1)\) and a convergence condition; finite horizons may use discount one.

\textbf{Optimization and existence.} A displayed \(\max\), \(\min\), or \(\arg\max\) is conditional on attainment. For finite feasible menus this is immediate; general spaces need continuity and compactness/coercivity or another existence argument. Without attainment, the target is the supremum/infimum and arbitrarily near-optimal feasible choices. \(\inf\varnothing=+\infty\). Empty feasibility and an unidentified target are different issues.

\textbf{Potential outcomes.} \(Y_i(z)\) is the outcome under a specified intervention, not the observed conditional mean among people choosing \(z\). In binary comparisons \(\Delta Y_i=Y_i(1)-Y_i(0)\). Specify assignment, treatment versions, compliance, information and observation horizon. Interference is allowed under a full policy vector; compressed exposure notation needs a justified mapping. No interference, consistency and overlap are substantive assumptions, not automatic consequences of notation.

\textbf{Populations and observation.} Fix baseline eligibility and population weights. Conditioning on post-policy employment, survival, relocation or program uptake can change the population being compared. Death is not ordinary loss to follow-up; outcomes truncated by death need a composite convention, joint survival target, or explicitly defined survivor stratum. Fertility-changing interventions need a population functional rather than an unexplained fixed descendant cohort. Measurement and missingness models must be supplied.

\textbf{Identification.} A structural model \(m\in\mathcal M\) implies observable law \(P_m\) and target \(\theta(m)\). Identification means
\[
 P_{m_1}=P_{m_2}\ \Longrightarrow\ \theta(m_1)=\theta(m_2)
 \quad\text{for all }m_1,m_2\in\mathcal M .
\]
Without identification, the target is
\[
 \Theta(P;\mathcal M)=\{\theta(m):m\in\mathcal M,\ P_m=P\}.
\]
Writing \(\theta(P)\) before establishing identification can hide incompatible latent models. Confidence sets additionally need a sampling law, confidence level, asymptotic sequence and uniformity class. Point identification, coverage and informative precision are separate properties.

\textbf{Mediation.} Total effects of randomized assignment are distinct from cross-world natural indirect effects. Belief/effort-channel effects require a meaningful mediator intervention and additional measurement, exclusion or mediation assumptions. Randomization of the initial policy alone does not supply those assumptions. Report bounds or interventional alternatives when the stronger target is unsupported.

\textbf{Equilibrium.} A counterfactual \(W(z)\), \(p(z)\), or \(\mu_t^z\) requires individual optimization, feasibility, government/resource budgets, market clearing and state-transition laws. State equilibrium existence and selection, or report ranges over compatible equilibria. Initial-state integration includes subsequent endogenous transitions; current-state integration must use the policy-dependent distribution. Reduced-form invariance after policy is not assumed.

\textbf{Welfare accounts.} Scores, utility, health, dollars and ecological quantities require explicit conversion before addition. Fees, taxes, subsidies, wages and extortion payments are transfers between parties; distributional weights can make them welfare relevant, but their face value is not automatically a resource loss. State ownership, geographic boundaries, foreign-income weights and fiscal finance. Do not add profits to household welfare already including dividends, or count land capitalization and the underlying benefit twice. A benefit-cost expression must distinguish gross benefits from net welfare and subtract every real cost exactly once.

\textbf{Derivatives and risk.} Log derivatives require positive arguments; local responses require admissible perturbations and specified equilibrium branches. Differentiation under expectations needs regularity/domination, and boundaries may allow only one-sided derivatives. Risk uses a specified law; ambiguity uses a declared nonempty law set on a common history space. Adaptive policies must be nonanticipative. A robust criterion is an editorial choice unless attributed explicitly.

\textbf{Scope overlaps.} Allocation, collective choice, contracts, household bargaining and coordinated policy overlap economics with optimization and game theory. They are retained because they appeared in the original catalogue; no claim of a strictly game-theory-free selection is made.
% END ECONOMICS STATEMENT
\end{document}
